\def\ChapterCode{NKA}
\chapter
[\CO{[NKA]} Notkompetenz: Arzneimittel]
{Notkompetenz: Arzneimittel}
\label{chp:NKA}




\ChapterIntro
{%

}
{%
\begin{description}
  \item[Maintainer:] Sebastian Gabriel
  \item[Autoren:] Diverse
  \item[Reviewer:] Standard-Reviewprozess
  \item[Version:] {\VarDocVersionTypeName} 
  ({\VarDocVersionTypeDescription})
  \item[SHA1:] {\DocGitCommit}
\end{description}

{\TxTeilkapitelAASS}
}

\section{\HeadingKeyword{Einleitung}}

\section{\HeadingKeyword{Verabreichungswege}}

\subsection{\HeadingKeyword{Intramuskuläre} Injektion}

\subsubsection{\HeadingKeyword{Allgemeines}}
Bei der intramuskulären Injektion ist eine
technisch und hygienisch richtige Anwendung besonders wichtig. Das
Risiko für folgende Komplikationen ist bei Nichtbeachten der u.\,g.
Handlungsanleitung erhöht:
 
\begin{compactenum}
\item Infektion
\item Nekrose
\item Abszess
\item Hämatom
\item Nervenschädigung
\end{compactenum}


\subsubsection{\HeadingKeyword{Material}}

\begin{FamPropMaterial}
\begin{compactenum}
  \item Nierentassen (rein/unrein)
  \item Untersuchungshandschuhe
  \item Einige saubere Tupfer
  \item Alkoholhältiges
  Hautantiseptikum
  \item Wundpflaster
  \item Spritze
  \item Medikament
  \item Kanüle zum
  Aufziehen des Medikaments (möglichst großes Lumen)
  \item Kanüle zur Injektion (möglichst kleines
  Lumen, Viskosität des Medikaments beachten)
  \item Ersatzinjektionskanüle
  \item Durchstichsicherer
  Entsorgungsbehälter (\enquote{Sharp-Box})
\end{compactenum}
\end{FamPropMaterial}




Die Kanülenlänge muss so bemessen werden, dass die
Kanülenspitze mindestens \EE{1\,cm} in der Muskulatur zu liegen kommt:

\E{z.\,B. Kleinkinder: \VonBis{30}{32}\MM*, normalgewichtige Erwachsene: 51\MM*}


\subsubsection{\HeadingKeyword{Handlungsschritte}}

\begin{FamPropHandlungsschritte}[{Vorbereitung und Kontaktaufnahme}]
\begin{enumerate}
  \item \StepHeading{Händedesinfektion}
  \item \StepHeading{Aufziehen des Injektionsmittels}
  \begin{enumerate}
    \item Siehe subkutane Injektion Punkt 2.
  \end{enumerate}
  \item \StepHeading{Verlassen des Stützpunktes} mit
  vorbereitetem Material
  \item \StepHeading{Kontaktaufnahme} mit dem Patienten
  \begin{enumerate}
    \item Kontrolle der Kontraindikationen: z.\,B. verminderte
    Blutgerinnung
  \end{enumerate}
  \item \StepHeading{Patientenlagerung}
  \begin{enumerate}
    \item Die Injektion wird am liegenden Patienten verabreicht
    \item Stabile Seitenlage
  \end{enumerate}
  \item \StepHeading{Handschuhe anziehen}
  \ListBreak[]{enumerate}{Punktion}
  \item \StepHeading{Bestimmen der Einstichstelle}
  \begin{enumerate}
    \item Beste Injektionsstelle: Hochstetter-Punkt (= ventrogluteale Injektion)
    \item Bei linker Seitenlage des Pat. (= Inj. rechts gluteal):
    Zeigefinger der linken Hand auf die Spina iliaca anterior superior,
    Mittelfinger maximal abgespreizt auf die Crista iliaca legen
    \item Die Injektion soll im unteren Drittel des so entstandenen Dreiecks
    erfolgen
%  
%\par\includegraphics[width=\linewidth]{./images/Skriptumaegf-img47.png}
  \end{enumerate}
  \item \StepHeading{Inspektion} der Einstichstelle
  \begin{enumerate}
    \item Kontraindikationen: Entzündungen, Schwellungen, Blutergüsse
    \item Die Injektionsstelle kann man sich durch Orientierung an Haaren,
    Muttermalen oder anderen Hautunreinheiten merken
    \item Niemals mit dem Fingernagel direkt die Injektionsstelle
    markieren (stattdessen wenn notwendig sterilisierte
    Sonde (Kanüle) wählen oder um Einstichstelle herum markieren)
  \end{enumerate}
  \item \StepHeading{Wischdesinfektion}
  \item \StepHeading{Stichschutz der Kanüle entfernen}
  \item \StepHeading{Injektion}
  \begin{enumerate}
    \item Spritze mit dominanter Hand wie Kugelschreiber fassen
    \item Der Einstich soll senkrecht zur Hautoberfläche und möglichst zügig
    durchgeführt werden
    \item Ein Sicherheitsabstand von 1\CM* vom Kanülenkonus zur Haut soll
    eingehalten werden
    \item Die Kanülenspitze muss im Muskelgewebe zu liegen kommen: d.\,h.
    tief genug einstechen, dass Kanülenspitze nicht im subkutanen
    Fettgewebe zu liegen kommt. Vorsicht: Das Periost des Hüftknochens
    ist sehr schmerzempfindlich.
    \item Aspirieren: Umgreifen, dazu die Spritze mit der nicht dominanten Hand
    fassen und fixieren, mit der dominanten Hand aspirieren
    \item Die Spritze \EE{um 90{\textdegree} drehen und
    nochmals aspirieren}\footnote{\E{90{\textdegree}-Drehung
    bei der Aspiration}: Rationale ist, dass sich in
    seltenen Fällen die Kanülenöffnung an der Gefäßwand
    ansaugen könnte, daher wird sicherheitshalber diese
    Drehung durchgeführt und ein zweites Mal aspiriert.
    Einige der i.m.-mäßig applizierten Substanzen können
    bei versehentlicher intravasaler Verabreichung sehr
    schwerwiegende Folgen haben.}; wenn auch dann kein Blut
    aspiriert wird, langsam injizieren. Wenn die
    Injektionslösung ölig ist, besonders langsam \item
    Patienten während der Injektion beobachten (Schmerzen?,
    Schwindel?)
  \end{enumerate}
  \item \StepHeading{Kanüle entfernen}
  \item \StepHeading{Kanüle entsorgen}
  \item \StepHeading{Tupfer fest auf die
  Einstichstelle drücken}, ca. 1 Minute
  \begin{enumerate}
    \item Wundpflaster
  \end{enumerate}
  \ListBreak[]{enumerate}{Nachbereitung und Verabschiedung}
  \item \StepHeading{Beenden des Patientenkontakts}
  \begin{enumerate}
    \item Bei Problemen sofort Arzt verständigen
  \end{enumerate}
  \item \StepHeading{Kontaminationsfreie Materialentsorgung}
  \item \StepHeading{Handschuhe ausziehen}
  \item \StepHeading{Händedesinfektion}
\end{enumerate}
\end{FamPropHandlungsschritte}





\subsubsection{Mögliche \HeadingKeyword{Probleme}}


\Text{{\SymbolKomplikationen} Blutaspiration}
{Injektion abbrechen und
Injektionsstelle einige Minuten fest komprimieren.}
{}
{}
{}
{}

\begin{Danger}
  \item Impfungen werden nicht intragluteal sondern in den M. deltoideus
  verabreicht.
\end{Danger}





\subsubsection{\RoteKarte\ \HeadingKeyword{Fallstricke}}
\begin{FamPropFallstricke}
\begin{itemize}
  \item[\RoteKarte] Selbst- oder Fremdgefährdung durch
  Kanülen, Nadeln (z.\,B. Entsorgung in den Restmüll oder Recapping)
  \item[\RoteKarte] Steriles Material wird kontaminiert und
  es erfolgt keine entsprechende Reaktion darauf
  \item[\RoteKarte] Aufziehkanüle wird am
  Nadelabwurfbehälter abgestreift und der Spritzenkonus kontaminiert
  \item[\RoteKarte] Fehlerhafte Hautantiseptik vor Punktion
  (z.\,B. vergessen, an falscher Stelle)
  \item[\RoteKarte] Fehlende oder mangelhafte
  Identitätsprüfung bei Verabreichung von Medikamenten
  \item[\RoteKarte] Patient wird mit Aufziehkanüle punktiert
\end{itemize}
\end{FamPropFallstricke}


\subsubsection{{\protect\dsliterary} Quellenangaben und
Weiterführende Literatur}
\cite{GabkaInjektionstechnik,InjektionenLeichtGemacht4,LippertAnatomie7}


\subsection{Sonderfall: 
\HeadingKeyword{Epipen{\texttrademark}}}


\section{Relevante \HeadingKeyword{Medikamente}}


%%%%%%%%%%%%%%%%%%%%%%%%%%%%%%%%%%%%%%%%%%%%%%%%%%%%%%%%%%%
%%%%%%%%%%%%%%%%%%%%%%%%%%%%%%%%%%%%%%%%%%%%%%%%%%%%%%%%%%%
\subsection{\HeadingKeyword{Katecholamine}}

%%%%%%%%%%%%%%%%%%%%%%%%%%%%%%%%%%%%%%%%%%%%%%%%%%%%%%%%%%%
%%%%%%%%%%%%%%%%%%%%%%%%%%%%%%%%%%%%%%%%%%%%%%%%%%%%%%%%%%%
\subsubsection{\HeadingKeyword{Adrenalin}}

\begin{labeling}[:]{mm}
  \item[Präp] \Conc{1}{10\,000}: L-Adrenalin, \Conc{1}{1\,000} Suprarenin.
  \item[Ind] Reanimation (schockbarer und nicht-schockbarer Rhythmus),
  Kardiogener Schock, Anaphylaxie > 3{\textdegree},  schweres Asthma bronchiale.
  \item[KI]Im Notfall keine, ansonsten:
  Hypertonie, KHK, Glaukom, Tachyarrhythmie, Hyperthyreose,
  Phäochromocytom.
  \item[\textrecipe]
  \begin{itemize}
    \item \EE{Reanimation:} Erw: 1\Mg, \EE{\E{Kinder}: 0,01\MgKg}
    (1\Ml*/\,10\,kg (bei 1:10\,000)); Wiederholung alle 3\,min.)
    \item \EE{Kardiogener Schock:}
    \VonBis{0,025}{0,3}\,{\textmu}g\,/\,kgKG\,/\,min
    \item \EE{Anaphylaxie:} 3{\textdegree} und \EE{schweres Asthma bronchiale}: 
    \E{Erw.}: 0,1\Mg* \Abk{IV} oder \VonBis{0,2}{0,3}\Mg* \Abk{SC}/\Abk{IM}; 
    \E{Kinder}: 0,01\MgKg/\Abk{SC}/\Abk{IM}; 
    Wiederholung alle \VonBis{5}{15} Min
    \item \EE{Kruppanfall:} \Conc{1}{1\,000} (sic!) im Vernebler \Conc{1}{2} mit \ce{NaCl}
    verdünnt, 0,1\MgKg* inhal.
  \end{itemize}
  \item[Anm] Bei tiefer endobronchialer Applikation \VonBis{2}{3}-fache Dosis.
  Bei gleichz. Ca-Gabe \EE{Arrhythmien}; \EE{NaBic} über gl. Zugang inakt.
  Katech.; Kein Adrenalin in Akren spritzen; sehr Temp.-labil
  ({\textgreater}\,40\,{\textcelsius}), nur farblose Lösungen verwenden; Je höher die
  Dosis, desto schlechter neurologisches Outcome.
  \item[Cave] Bei Betablocker-Pat. überwiegend
  {$\alpha$}-Rezeptoren-Wirkung.
\end{labeling}

\subsection{\HeadingKeyword{Narkose}}





%%%%%%%%%%%%%%%%%%%%%%%%%%%%%%%%%%%%%%%%%%%%%%%%%%%%%%%%%%%
%%%%%%%%%%%%%%%%%%%%%%%%%%%%%%%%%%%%%%%%%%%%%%%%%%%%%%%%%%%
\subsubsection{\HeadingKeyword{Diazepam}}
\begin{labeling}[:]{mm}
  \item[Präp] Amp: Gewacalm (2\Ml*/\,10\Mg*), Gtt: Psychopax Drops,
  Rectal: 5\Mg, 10\Mg
  \item[Beschreibung]
  \item[Ind] Sedierung, Epilepsie, Eklampsie, Prämediaktion (CV), Rectal bei
  Krämpfen
  \item[KI] chron. Hyperkapnie, Myasthenia gravis, Intox mit
zentral dämpfenden Substanzen
\item[Rx i.V.]  0,2-0,5\MgKg*, maximal 20\Mg
\item[Cave] CAVE: Kumulation. Nicht bei Kindern {\textless}2\,a, auch
imAtemdepression
\item[Rx rect.] $\leq$ 15\,kg: 5\Mg, {\textgreater}15\,kg 10\Mg
\item[Anm.] Antidot: Flumazenil (Anexate)

\end{labeling}


%%%%%%%%%%%%%%%%%%%%%%%%%%%%%%%%%%%%%%%%%%%%%%%%%%%%%%%%%%%
%%%%%%%%%%%%%%%%%%%%%%%%%%%%%%%%%%%%%%%%%%%%%%%%%%%%%%%%%%%
\subsubsection{\HeadingKeyword{Flumazenil}}
\opt{STATUS-EXPERIMENTAL}{

\begin{labeling}[:]{mm}
  \item[Präp] Anexate (5\Ml*/\,0,5mg, 10\Ml*/\,1\Mg*)
  \item[Beschreibung] Benzoantagonist
  \item[Ind] \Sign{↑} ICP, Epi, Intox mit Trizyklika
  \item[KI]
  \item[Rx.] initial \VonBis{0,2}{0,3}\Mg, Wiederholung 0,2\Mg* alle 60\,sek,
  Gesamtdosis: \VonBis{1}{2}\Mg. \E{Perfusor}: \VonBis{0,1}{0,4}\Mg/h;
  \E{Kinder}: \VonBis{4}{20}\,\nicefrac{{\textmu}g}{kg}
  \item[Anm.] evt. Auslösung von Entzugssymptomen, Temp.-stabil
\end{labeling}

} % STATUS-EXPERIMENTAL




%%%%%%%%%%%%%%%%%%%%%%%%%%%%%%%%%%%%%%%%%%%%%%%%%%%%%%%%%%%
%%%%%%%%%%%%%%%%%%%%%%%%%%%%%%%%%%%%%%%%%%%%%%%%%%%%%%%%%%%

\subsubsection{\HeadingKeyword{Midazolam}}
{
\begin{labeling}[:]{mm}
  \item[Präp] Dormicum (1\Ml*/\,5mg, 3\Ml*/\,15mg, 5\Ml*/\,5mg (!!), 10\Ml*/\,50\Mg*)
  \item[Beschreibung]
  \item[Ind] Sedierung, zerebraler Krampfanfall,
  Status epi, Narkoseeinleitungchron, Hyperkapnie, Myasthenia gravis,
  ImzdS
  \item[KI]
  \item[Rx] Sedierung: \VonBis{0,015}{0,03} \MgKg\newline
  Krampf: \VonBis{0,05}{0,1}\MgKg\newline
  Status: bis 0,2\MgKg\newline
  Einleitung: \VonBis{0,1}{0,2}\MgKg
  \item[Anm.] Atemdepression gut steuerbar, Wirkdauder individuell
  \item[Cave]  paradoxe Wirkung va bei älteren
\end{labeling}
}


%%%%%%%%%%%%%%%%%%%%%%%%%%%%%%%%%%%%%%%%%%%%%%%%%%%%%%%%%%%
%%%%%%%%%%%%%%%%%%%%%%%%%%%%%%%%%%%%%%%%%%%%%%%%%%%%%%%%%%%
\subsubsection{Opiate}
%%%%%%%%%%%%%%%%%%%%%%%%%%%%%%%%%%%%%%%%%%%%%%%%%%%%%%%%%%%
%%%%%%%%%%%%%%%%%%%%%%%%%%%%%%%%%%%%%%%%%%%%%%%%%%%%%%%%%%%
\subsubsection{Fentanyl}


\begin{labeling}[:]{mm}
  \item[Präp] (2\Ml*/\,0,1mg=100{\textmu}g, 10\Ml*/\,0,5\Mg*)
  \item[Beschreibung:]
  \item[Ind] Opioidanalg., Einl. und Vert. d. Narkose
  \item[KI] Im Notfall keine, sonst: Stillen, Sgl., Asthma, akuter hepatischer
  Porphyrie, Myasthenia gravis, Geburtshilfe (Sektio)
  \item[Rx] Erw: \VonBis{0,05}{0,2}\Mg; \E{Kinder}: \VonBis{1}{2}\,(--\,5)
  {\textmu}g/kg \newline Dosisred. bei LPS, NINS, aälteren, schlehctem AZ.
  \item[Anm.] Tip: 0,1\,\Ml*/\,5\,{\textmu}g. {\textmu}-Agonist, Atemdepress.,
  insbes. bei Komb. mit Benzos. In höheren Dosen Thoraxrigidität. Antagonist:
  Naloxon.
\end{labeling}




%%%%%%%%%%%%%%%%%%%%%%%%%%%%%%%%%%%%%%%%%%%%%%%%%%%%%%%%%%%
%%%%%%%%%%%%%%%%%%%%%%%%%%%%%%%%%%%%%%%%%%%%%%%%%%%%%%%%%%%
\subsubsection{Morphin (hydrochlorid)}

Vendal (1\Ml*/\,10\Mg*)

\begin{labeling}[:]{mm}
  \item[Präp] ---.
  \item[Beschreibung] ---.
  \item[Ind] Analgesie insbes. Herzinfarkt.
  \item[KI]  im Notfall: Asthma, spastische Schmerzen, Sgl, Stillen.
  \item[Rx] Erw: \VonBis{2,5}{10}\Mg* \Abk{IV}, \E{Kinder}: 0,05\MgKg* iv; auch
  sc/im.
  \item[Cave] ---.
  \item[Anm] {\textmu}-Agonist, Atemdepression, parasympmimetische NW.
\end{labeling}






%%%%%%%%%%%%%%%%%%%%%%%%%%%%%%%%%%%%%%%%%%%%%%%%%%%%%%%%%%%
%%%%%%%%%%%%%%%%%%%%%%%%%%%%%%%%%%%%%%%%%%%%%%%%%%%%%%%%%%%
\subsubsection{Naloxon}


\begin{labeling}[:]{mm}
  \item[Präp]  Narcanti (1\Ml*/\,0,4\Mg*).
  \item[Beschreibung] Opiat-Antagonist .
  \item[Ind] ---.
  \item[KI] Überempfiundlichkeit, Atemdepression die nicht durch Opiate
ausgelöst wurde.
  \item[Rx] Erw: 0,4\Mg* \Abk{IV}, \Abk{IM}, sc\newline
\E{Kinder}: 0,1\MgKg* \Abk{IV}, \Abk{IM}, sc.
  \item[Cave]  Cave bei Pat mit vorbestehender Herzerkr. .
  \item[Anm]  {\textmu}-Antagonist, daher Entzugssympt., auch bei Gebärenden
indiziert---.
\end{labeling}


\subsection{Atmung}

%%%%%%%%%%%%%%%%%%%%%%%%%%%%%%%%%%%%%%%%%%%%%%%%%%%%%%%%%%%
%%%%%%%%%%%%%%%%%%%%%%%%%%%%%%%%%%%%%%%%%%%%%%%%%%%%%%%%%%%
\subsubsection{Fenoterol}

\begin{labeling}[:]{mm}
  \item[Präp] \Z{
  Berotec Dae (0,2mg/Hub); Partusisten (1\Ml*/\,0,025mg, 1\Ml*/\,0,5\Mg*)}.
  \item[Beschreibung] ---.
  \item[Ind] ---.
  \item[KI]  Dae: im Notfall keine. Tokolyse \Abk{IV}: starke Blutung, vorz
  Plazentalösung, Tachyarrhythmie, Thyreotoxikose, hypertrophe obstruktive
  CMP---.
  \item[Rx]
  \begin{inparaitem}
  \item Asthma Bronchiale: \VonBis{1}{2} Hübe, Wiederholung nach
  \VonBis{5}{10}{\textquotesingle}
  \item Tokolyse mit Dosieraerosol: \VonBis{2}{5} Hübe
  \item Tokolyse \Abk{IV}: sehr langsamer Bolus: 25\,{\textmu}g.
  \E{Perfusor}: \VonBis{0,01}{0,05}\,{\textmu}g\,/\,kgKG\,/\,min., zB
  2\Mg\,/\,50\Ml*/\,1-5 
  \end{inparaitem}.
  \item[Cave] ---.
  \item[Anm] DAe: Im Anfall meist unwirksam. \Abk{IV}: Kontrolle von Herzfrequenz von Mutter u
  Kind. Optional Hexoprenalin (Gynipral)..
\end{labeling}


\subsubsection{Azetylsalizylsäure (ASS)}

\subsubsection{Nitroglycerin}



\subsection{Infusionslösungen -- allgemein}

\subsection{Kristalloide Infusionslösungen}




\section{Algorithmen}
\ChapterEnd